% Options for packages loaded elsewhere
\PassOptionsToPackage{unicode}{hyperref}
\PassOptionsToPackage{hyphens}{url}
%
\documentclass[
]{article}
\usepackage{lmodern}
\usepackage{amssymb,amsmath}
\usepackage{ifxetex,ifluatex}
\ifnum 0\ifxetex 1\fi\ifluatex 1\fi=0 % if pdftex
  \usepackage[T1]{fontenc}
  \usepackage[utf8]{inputenc}
  \usepackage{textcomp} % provide euro and other symbols
\else % if luatex or xetex
  \usepackage{unicode-math}
  \defaultfontfeatures{Scale=MatchLowercase}
  \defaultfontfeatures[\rmfamily]{Ligatures=TeX,Scale=1}
\fi
% Use upquote if available, for straight quotes in verbatim environments
\IfFileExists{upquote.sty}{\usepackage{upquote}}{}
\IfFileExists{microtype.sty}{% use microtype if available
  \usepackage[]{microtype}
  \UseMicrotypeSet[protrusion]{basicmath} % disable protrusion for tt fonts
}{}
\makeatletter
\@ifundefined{KOMAClassName}{% if non-KOMA class
  \IfFileExists{parskip.sty}{%
    \usepackage{parskip}
  }{% else
    \setlength{\parindent}{0pt}
    \setlength{\parskip}{6pt plus 2pt minus 1pt}}
}{% if KOMA class
  \KOMAoptions{parskip=half}}
\makeatother
\usepackage{xcolor}
\IfFileExists{xurl.sty}{\usepackage{xurl}}{} % add URL line breaks if available
\IfFileExists{bookmark.sty}{\usepackage{bookmark}}{\usepackage{hyperref}}
\hypersetup{
  hidelinks,
  pdfcreator={LaTeX via pandoc}}
\urlstyle{same} % disable monospaced font for URLs
\usepackage{longtable,booktabs}
% Correct order of tables after \paragraph or \subparagraph
\usepackage{etoolbox}
\makeatletter
\patchcmd\longtable{\par}{\if@noskipsec\mbox{}\fi\par}{}{}
\makeatother
% Allow footnotes in longtable head/foot
\IfFileExists{footnotehyper.sty}{\usepackage{footnotehyper}}{\usepackage{footnote}}
\makesavenoteenv{longtable}
\setlength{\emergencystretch}{3em} % prevent overfull lines
\providecommand{\tightlist}{%
  \setlength{\itemsep}{0pt}\setlength{\parskip}{0pt}}
\setcounter{secnumdepth}{-\maxdimen} % remove section numbering

\author{}
\date{}

\begin{document}

\hypertarget{internal-benchmark-methodology-for-fungus-sv}{%
\section{Internal Benchmark Methodology for
FUNGUS-SV}\label{internal-benchmark-methodology-for-fungus-sv}}

\hypertarget{assembly-based-truth-set-construction-for-non-model-organisms}{%
\subsection{Assembly-Based Truth Set Construction for Non-Model
Organisms}\label{assembly-based-truth-set-construction-for-non-model-organisms}}

\hypertarget{authors-kelton-jenkov-guimaruxe3es-et-al.}{%
\subsubsection{Authors: Kelton Jenkov Guimarães et
al.}\label{authors-kelton-jenkov-guimaruxe3es-et-al.}}

\hypertarget{pipeline-fungus-sv-v0.5.1}{%
\subsubsection{Pipeline: FUNGUS-SV
v0.5.1}\label{pipeline-fungus-sv-v0.5.1}}

\hypertarget{date-may-2026}{%
\subsubsection{Date: May 2026}\label{date-may-2026}}

\begin{center}\rule{0.5\linewidth}{0.5pt}\end{center}

\hypertarget{background}{%
\subsection{1. Background}\label{background}}

Structural variant (SV) detection in non-model organisms lacks the
gold-standard benchmark sets available for human genomes (GIAB, HGSVC,
SMaHT). Without a truth set, it is impossible to measure precision,
recall, or false discovery rates for any SV caller. This document
describes a computational methodology for constructing an internal truth
set using de novo genome assembly, enabling rigorous validation of SV
detection pipelines in any organism.

The methodology is based on approaches used by: - \textbf{Kronenberg et
al.~(2025)} \emph{Nature Methods} --- Platinum Pedigree truth set
construction - \textbf{Zhang et al.~(2025)} \emph{bioRxiv} --- SMaHT
mosaic SV benchmark - \textbf{Ebert et al.~(2021)} \emph{Science} ---
HGSVC assembly-based SV discovery - \textbf{Chen et al.~(2023)}
\emph{Nature Communications} --- DeBreak assembly validation

\begin{center}\rule{0.5\linewidth}{0.5pt}\end{center}

\hypertarget{principle}{%
\subsection{2. Principle}\label{principle}}

The core principle is: \textbf{SVs detected by comparing a de novo
assembly to a reference genome are independent of read-alignment
artifacts that affect read-based SV callers.} When both an
assembly-based approach and a read-based approach (FUNGUS-SV) agree on
an SV, the probability that it is a true variant is substantially higher
than either method alone.

\hypertarget{why-assembly-based-svs-are-more-reliable}{%
\subsubsection{Why Assembly-Based SVs Are More
Reliable}\label{why-assembly-based-svs-are-more-reliable}}

\begin{longtable}[]{@{}lll@{}}
\toprule
\begin{minipage}[b]{0.39\columnwidth}\raggedright
Evidence Type\strut
\end{minipage} & \begin{minipage}[b]{0.22\columnwidth}\raggedright
Source\strut
\end{minipage} & \begin{minipage}[b]{0.30\columnwidth}\raggedright
Artifacts\strut
\end{minipage}\tabularnewline
\midrule
\endhead
\begin{minipage}[t]{0.39\columnwidth}\raggedright
Read-alignment SVs\strut
\end{minipage} & \begin{minipage}[t]{0.22\columnwidth}\raggedright
Mapped reads (BAM)\strut
\end{minipage} & \begin{minipage}[t]{0.30\columnwidth}\raggedright
Mapping errors, chimeric alignments, repeat collapse\strut
\end{minipage}\tabularnewline
\begin{minipage}[t]{0.39\columnwidth}\raggedright
Assembly-based SVs\strut
\end{minipage} & \begin{minipage}[t]{0.22\columnwidth}\raggedright
Contigs from de novo assembly\strut
\end{minipage} & \begin{minipage}[t]{0.30\columnwidth}\raggedright
Assembly errors (less frequent with HiFi)\strut
\end{minipage}\tabularnewline
\begin{minipage}[t]{0.39\columnwidth}\raggedright
\textbf{Concordant SVs}\strut
\end{minipage} & \begin{minipage}[t]{0.22\columnwidth}\raggedright
\textbf{Both agree}\strut
\end{minipage} & \begin{minipage}[t]{0.30\columnwidth}\raggedright
\textbf{Minimal --- orthogonal methods confirm each other}\strut
\end{minipage}\tabularnewline
\bottomrule
\end{longtable}

The assembly-based approach has been validated in human genomics, where
assembly-based SV calls from hifiasm + PAV achieve \textgreater90\%
concordance with GIAB truth sets (Ebert et al.~2021).

\begin{center}\rule{0.5\linewidth}{0.5pt}\end{center}

\hypertarget{methodology}{%
\subsection{3. Methodology}\label{methodology}}

\hypertarget{required-data}{%
\subsubsection{3.1 Required Data}\label{required-data}}

\begin{longtable}[]{@{}lll@{}}
\toprule
Item & Format & Purpose\tabularnewline
\midrule
\endhead
PacBio HiFi reads & FASTQ (.fastq.gz) & From the isolate to be
analyzed\tabularnewline
Reference genome & FASTA (.fasta) & Different strain/species for
comparison\tabularnewline
\bottomrule
\end{longtable}

\hypertarget{step-1-de-novo-genome-assembly}{%
\subsubsection{3.2 Step 1: De Novo Genome
Assembly}\label{step-1-de-novo-genome-assembly}}

Assemble the HiFi reads into a haplotype-resolved genome assembly using
hifiasm:

```bash \# Install hifiasm conda create -n assembly -c bioconda hifiasm
-y conda activate assembly

\hypertarget{run-assembly-haploid-mode-for-fungi}{%
\section{Run assembly (haploid mode for
fungi)}\label{run-assembly-haploid-mode-for-fungi}}

hifiasm -o isolate\_asm -t 16 --primary isolate\_reads.fastq.gz

\hypertarget{convert-to-fasta}{%
\section{Convert to FASTA}\label{convert-to-fasta}}

awk `/\^{}S/\{print ``\textgreater{}''\$2;print \$3\}'
isolate\_asm.p\_ctg.gfa \textgreater{} isolate\_asm.p\_ctg.fasta
Expected output: A haploid assembly with contig N50 \textgreater{} 1 Mb
for fungal genomes at \textgreater50× HiFi coverage.

Quality control:

\hypertarget{check-assembly-statistics}{%
\section{Check assembly statistics}\label{check-assembly-statistics}}

quast isolate\_asm.p\_ctg.fasta -r reference.fasta -o quast\_results/
3.3 Step 2: Assembly-Based SV Calling Compare the assembled genome to
the reference using SVIM-asm in haploid mode: \# Align assembly to
reference minimap2 -x asm5 -t 8 --cs reference.fasta
isolate\_asm.p\_ctg.fasta\\
\textbar{} samtools sort -@ 4 -o asm\_to\_ref.sorted.bam - samtools
index asm\_to\_ref.sorted.bam

\hypertarget{call-svs-from-assembly-comparison-haploid-mode}{%
\section{Call SVs from assembly comparison (haploid
mode)}\label{call-svs-from-assembly-comparison-haploid-mode}}

svim-asm haploid --min\_sv\_size 50 asm\_to\_ref.sorted.bam
reference.fasta\\
results/asm\_svs/ Why SVIM-asm:

Designed for both haploid and diploid assemblies (Heller \& Vingron
2020)

Validated on human HGSVC data with high precision

Reports DEL, INS, INV, DUP, and complex SVs

3.4 Step 3: Read-Based SV Detection (FUNGUS-SV) Run FUNGUS-SV on the
same HiFi reads aligned to the same reference:

\hypertarget{align-reads-to-reference}{%
\section{Align reads to reference}\label{align-reads-to-reference}}

minimap2 -t 8 -ax map-hifi -R `@RG\tID:isolate\tSM:sample'\\
reference.fasta isolate\_reads.fastq.gz\\
\textbar{} samtools sort -@ 4 -o reads\_to\_ref.sorted.bam - samtools
index reads\_to\_ref.sorted.bam

\hypertarget{run-fungus-sv-icb-consensus}{%
\section{Run FUNGUS-SV ICB
consensus}\label{run-fungus-sv-icb-consensus}}

python fungus\_sv/core/icb.py\\
--bam reads\_to\_ref.sorted.bam\\
--reference reference.fasta\\
--output results/icb/\\
--callers sniffles2 cutesv svim\\
--min-callers 2

\hypertarget{run-fungus-sv-validation}{%
\section{Run FUNGUS-SV validation}\label{run-fungus-sv-validation}}

python -m valid\_sv.run\_validation\\
--consensus-vcf results/icb/consensus\_svs.vcf\\
--bam reads\_to\_ref.sorted.bam\\
--reference reference.fasta\\
--fastq isolate\_reads.fastq.gz\\
--output results/validation/ 3.5 Step 4: Benchmark Comparison Compare
FUNGUS-SV results to the assembly-based truth set using Truvari: \#
Compare ICB consensus to assembly SVs truvari bench\\
-b results/asm\_svs/variants.vcf\\
-c results/icb/consensus\_svs.vcf\\
--reference reference.fasta\\
-o benchmark\_results/\\
--pctseq 0 --pctsize 0.7 --refdist 500

\hypertarget{generate-summary-metrics}{%
\section{Generate summary metrics}\label{generate-summary-metrics}}

truvari summarize --input benchmark\_results/ --output summary.txt 3.6
Step 5: Stratified Performance Analysis Compute precision, recall, and
F1 score stratified by:

Stratification Why SV type (DEL, INS, INV, DUP) Different types have
different detection difficulty SV size (50-100 bp, 100-500 bp, 500-5000
bp, \textgreater5 kb) Size-dependent accuracy T-score tier (HIGH ≥0.6,
MEDIUM 0.4-0.6, WEAK \textless0.4) Validates T-score calibration Repeat
context (in repeat vs.~unique) Repeat regions are challenging Number of
supporting ICB callers (2 vs.~3) Validates ICB consensus 4. Expected
Outcomes 4.1 Performance Targets (Based on Human Benchmarks) Metric
Target Source ICB Precision \textgreater80\% SV-JIM (Todd et al.~2025):
3-caller consensus ICB Recall (DEL) \textgreater60\% Liu et al.~(2024):
best callers on PacBio ICB Recall (INS) \textgreater40\% Liu et
al.~(2024): insertions are harder T≥0.6 FDR \textless15\% SV-MeCa
(Nkouamedjo et al.~2025) Ablation: multi-layer \textgreater{}
single-layer ≥15\% improvement Expected from triangulation theory 4.2
T-Score Calibration The benchmark enables mapping T-score thresholds to
empirical false discovery rates:

T-Score Expected FDR Action ≥0.8 \textless5\% Suitable for functional
follow-up ≥0.6 \textless15\% Prioritized candidate ≥0.4 \textless35\%
Exploratory only \textless0.4 \textgreater50\% Likely false positive
These thresholds are estimates pending empirical calibration.

\begin{enumerate}
\def\labelenumi{\arabic{enumi}.}
\setcounter{enumi}{4}
\tightlist
\item
  Limitations Assembly errors can create false SVs. hifiasm assemblies
  of HiFi data have high accuracy (QV \textgreater{} 50), but errors in
  repetitive regions can introduce false SVs that will be misclassified
  as true positives.
\end{enumerate}

Assembly-based callers miss complex SVs. SVIM-asm does not detect
translocations or complex rearrangements. The truth set will be
incomplete.

This is a silver standard, not gold. Unlike GIAB (which uses multiple
technologies, pedigree information, and manual curation), this approach
uses a single technology and automated methods.

Validated only for haploid genomes. Diploid or polyploid genomes require
phasing before assembly, which introduces additional complexity.

\begin{enumerate}
\def\labelenumi{\arabic{enumi}.}
\setcounter{enumi}{5}
\tightlist
\item
  Validation Against Published Methods To ensure the benchmark is
  reliable, cross-validate against orthogonal approaches:
\end{enumerate}

Method Tool Comparison Different assembly-based caller PAV Compare to
SVIM-asm results Different assembler Flye Compare to hifiasm assembly
Read-only validation SVvalidation (Zheng 2024) Independent breakpoint
confirmation High concordance (\textgreater80\%) between methods
increases confidence in the truth set.

\begin{enumerate}
\def\labelenumi{\arabic{enumi}.}
\setcounter{enumi}{6}
\tightlist
\item
  Software Requirements Tool Version Installation hifiasm ≥0.19 conda
  install -c bioconda hifiasm minimap2 ≥2.30 conda install -c bioconda
  minimap2 samtools ≥1.17 conda install -c bioconda samtools svim-asm
  ≥1.0 conda install -c bioconda svim-asm Truvari ≥5.0 conda install -c
  bioconda truvari quast ≥5.0 conda install -c bioconda quast
\item
  References Kronenberg, Z. et al.~(2025). The Platinum Pedigree: A
  long-read benchmark for genetic variants. Nature Methods.
\end{enumerate}

Zhang, Y. et al.~(2025). Comprehensive benchmarking of somatic
structural variant detection at ultra-low allele fractions. bioRxiv.

Ebert, P. et al.~(2021). Haplotype-resolved diverse human genomes and
integrated analysis of structural variation. Science, 372.

Chen, Y. et al.~(2023). Deciphering the exact breakpoints of structural
variations using long sequencing reads with DeBreak. Nature
Communications, 14.

Heller, D. \& Vingron, M. (2020). SVIM-asm: structural variant detection
from haploid and diploid genome assemblies. Bioinformatics, 36.

Liu, Y.H. et al.~(2024). Tradeoffs in alignment and assembly-based
methods for structural variant detection with long-read sequencing data.
Nature Communications, 15.

Todd, C. et al.~(2025). SV-JIM: detailed pairwise structural variant
calling using long-reads and genome assemblies. Methods, 234.

Nkouamedjo Fankep, R.C. et al.~(2025). SV-MeCa: an XGBoost-based
meta-caller approach for structural variant calling. BMC Bioinformatics,
26.

Zheng, Y. \& Shang, X. (2024). SVvalidation: A long-read-based
validation method for genomic structural variation. PLOS ONE, 19.

This methodology is part of the FUNGUS-SV pipeline documentation. The
pipeline is available at
https://github.com/keltonjenkovguimaraes-alt/fungus-sv

\end{document}
